\honest{The Fallacy of Gray}{Eliezer Yudkowsky}{2008}

I open with a dialogue from Marc Stiegler's novel \textsc{David's Sling}. The Sophisticate says that the world is all gray, so no one is better than anyone else. The Zetet replies that the Sophisticate has swapped a two-color view for a one-color view.

I do not know whether this mistake has an official name, so I call it the Fallacy of Gray. \nb{It has two names already: the nirvana fallacy, and false equivalence.} My example is a man from my previous post, ``But There's Still A Chance, Right?'' He was told that the odds against him were two to the power of 750 million to one, and said there was still a chance. To him all probabilities were simply ``uncertain,'' so he could ignore them.

That the Moon is made of green cheese and that the Sun is made mostly of hydrogen and helium are both uncertain, but not equally. Everything is gray, but some grays are nearly white. Or, if not, we can at least say which of two grays is darker.

A paragraph from Iain M. Banks's novel \textsc{The Player of Games} taught me this, and I set one clause in bold: some societies try to maximize the values they impose on their members, and some try to minimize them. The rest of the paragraph says that neutrality is probably impossible and that your politics cannot be separated from the rest of you.

The idea, I admit, is ``so obvious in retrospect that I could have conceivably picked it up in a hundred places.'' It is this: what you cannot switch off, you can still turn down. That is the whole essay.

Is it too obvious to mention? No, because ``many bloggers'' have said that nobody can completely eliminate bias. I name none of them.

Then comes my one named opponent, the economist Tyler Cowen. Robin Hanson, who writes this blog with me, said he gives at least 75\% weight to economic theory over his intuitions and tries to apply it straightforwardly. Cowen replied that theories are always applied through personal and cultural filters. But you can try to minimize that effect, and then call the result ``straightforward.'' \nb{Cowen's point was that economic theory has biases of its own, and that judgment can offset them. That is a claim about which gray is lighter, which is the essay's own subject. The post uses him as an example of treating all grays as the same.}

``Everyone is imperfect,'' but Gandhi and Stalin were not the same shade of imperfect. Saying ``some people are less imperfect than others'' wins no applause, but it gives hope to those who try. When told perfectionism is bad for me, I answer that it is fine to be imperfect, ``but not so imperfect that other people notice.''

People who say that every scientific paradigm shapes how it reads experiments act as if science were no better than witchdoctoring. Every worldview imposes some of its structure on what it sees, but some try to minimize that and some glory in it. I name none of them.

Expecting the Moon tomorrow because it has always been there, and expecting an invisible dragon to cure your daughter, are different degrees of uncertainty: one expects what has happened, to twelve decimal places, to happen again; the other expects something that breaks the observed order. Calling both ``faith'' is too broad. \nb{That is true. But whether you may assume that the future will resemble the past is exactly what the objection ``Science is based on faith too'' is about, and the post simply assumes it. Six months later, in ``Where Recursive Justification Hits Bottom,'' Yudkowsky wrote that a defense of induction, pressed far enough, goes ``around in a loop'': one expects the future to resemble the past because that rule has usually worked before.}

People who call faith good say ``Science is based on faith too!'' in an angry, triumphant tone, not as a compliment. Then I address the ``faithist'' directly: if science is a religion, it is the one that heals the sick and walks on the Moon. \nb{That answers the comparison, not the claim. The claim was about how science justifies its starting assumptions, and walking on the Moon relies on the same induction.}

Criticism should name ``specific problems X, Y, and Z.'' People who say ``It's gray!'' without naming any may have ``made a devil's bargain with their own mistakes'': having found an excuse not to improve, they deny that anyone can, and say proudly ``I'm glad to be gray'' and angrily ``And you're gray too!'' I name none of these people either.

In my last line I credit a commenter for pointing to Isaac Asimov's ``The Relativity of Wrong'': the flat Earth and the spherical Earth were both wrong, but whoever thinks them equally wrong is ``wronger than both of them put together.'' \nb{That essay had made the central point, for scientific theories, about twenty years earlier.}
